% ---
% INFORMAÇÕES REFERENTE A DISSERTAÇÃO
% ---


\titulo{Comunicação síncrona via REST e assíncrona orientada a eventos: uma análise experimental no processamento de registros}

% Autores
\autor{Ilan Wendling Thoele}
\autorcitado{Thoele, Ilan Wendling}
\segundoautor{}
\segundoautorcitado{}
\terceiroautor{}
\terceiroautorcitado{}

% Orientador e Coorientador
\orientador{Alessandra Bussador}
\orientadortitulo{Professora}
\coorientador{}
\coorientadortitulo{}

% Nomes dos membros da banca
\profbancaum{}
\profbancaumtitulo{}
\profbancadois{}
\profbancadoistitulo{}
\instituicaoprofum{}
\instituicaoprofdois{}

% Ano de realização do trabalho
\data{2026}
% Data de realização da banca de avaliação do trabalho
\databanca{} % Preencher somente após confirmação institucional; folha não emitida.

% Informações para a ficha catalográfica
\codigoautor{} % Código a ser informado pela biblioteca; ficha não emitida.
\codigopublicacao{}

% Informações de caracterização do trabalho
\palavrachaveum{Arquitetura orientada a eventos}
\palavrachavedois{REST}
\palavrachavetres{Registros de auditoria}
\palavrachavequatro{Apache Kafka}
\textoresumo{Este trabalho propõe comparar experimentalmente duas formas de comunicação aplicadas ao processamento de registros de auditoria: uma implementação síncrona baseada em API REST e uma implementação assíncrona orientada a eventos com Apache Kafka. As variantes Java compartilham o mesmo contrato de dados, a mesma regra de validação e a mesma persistência em PostgreSQL. A comparação futura considerará métricas de desempenho e cenários de indisponibilidade e recuperação, sem pressupor superioridade de uma abordagem. A implementação contempla confirmação de publicação, persistência idempotente, tentativas limitadas e retenção de falhas terminais. Testes unitários e verificações funcionais de integração contínua examinam os caminhos principais e a interrupção do banco. A pesquisa documenta a fundamentação, o método e o estado do protótipo Java; a instrumentação quantitativa, o piloto, o congelamento do protocolo e as repetições experimentais permanecem pendentes.}
\keywordum{Event-driven architecture}
\keyworddois{REST}
\keywordtres{Audit records}
\keywordquatro{Apache Kafka}
\textoabstract{This work proposes an experimental comparison between two communication approaches for audit record processing: a synchronous implementation based on a REST API and an asynchronous event-driven implementation using Apache Kafka. Both variants share the same data contract, validation rules, and PostgreSQL persistence. Future comparison will address performance metrics and failure and recovery scenarios without assuming the superiority of either approach. This version documents the theoretical background, planned method, and current prototype status, without reporting experimental results.}
\textodedicatoria{}
\textoepigrafe{}
\textoagradecimento{}


% ---
% INFORMAÇÕES DO CURSO E INSTITUIÇÃO
% ---

\local{Foz do Iguaçu}
\instituicao{Centro Universitário Dinâmica das Cataratas}
\instituicaoSigla{UDC}
\tipotrabalho{Trabalho de Conclusão de Curso}
\curso{Sistemas de Informação}
\preambulo{\imprimirtipotrabalho\ apresentado como requisito obrigatório para obtenção do título de Bacharel em \imprimircurso\ do \imprimirinstituicao.}
\textoaprovacao{\imprimirtipotrabalho\ aprovado como requisito obrigatório para obtenção do título de Bacharel em \imprimircurso\ do \imprimirinstituicao, pela seguinte banca examinadora:}
